% ACI Companion — Lesson 1: AI Agents, PEAS & Environments
\documentclass[a4paper,11pt]{article}
\usepackage[T1]{fontenc}
\usepackage[utf8]{inputenc}
\usepackage{lmodern}
\usepackage{microtype}
\usepackage[a4paper,margin=1in,headheight=15pt,headsep=14pt]{geometry}
\usepackage{amsmath,amssymb}
\usepackage{graphicx}
\usepackage{booktabs}
\usepackage[table,svgnames]{xcolor}
\usepackage[most]{tcolorbox}
\tcbuselibrary{skins,breakable}
\usepackage{pifont}
\usepackage{enumitem}
\setlist{leftmargin=1.5em,topsep=4pt,itemsep=3pt}
\usepackage{parskip}
\usepackage{fancyhdr}
\usepackage{titlesec}
\usepackage[hidelinks]{hyperref}
\usepackage{xurl}

% Colours
\definecolor{NavyBanner}{HTML}{21355E}
\definecolor{IntFrame}{HTML}{2C5AA0}
\definecolor{IntTint}{HTML}{EBF0FA}
\definecolor{EveFrame}{HTML}{8A5A1E}
\definecolor{EveTint}{HTML}{FDF3E7}
\definecolor{WrkFrame}{HTML}{2E7D52}
\definecolor{WrkTint}{HTML}{EBF5EF}
\definecolor{WatFrame}{HTML}{B23A48}
\definecolor{WatTint}{HTML}{FAE9EB}
\definecolor{KeyFrame}{HTML}{6A4C93}
\definecolor{KeyTint}{HTML}{F0EBF8}
\definecolor{SectionNavy}{HTML}{1A2744}

% Callout boxes
\newtcolorbox{intuition}{enhanced,breakable,colback=IntTint,colframe=IntFrame,
  boxrule=0.6pt,arc=4pt,
  attach boxed title to top left={xshift=6mm,yshift=-3mm},
  boxed title style={colback=IntFrame,rounded corners},
  coltitle=white,fonttitle=\bfseries\sffamily\small,
  title={\ding{43}\; The intuition}}
\newtcolorbox{everyday}{enhanced,breakable,colback=EveTint,colframe=EveFrame,
  boxrule=0.6pt,arc=4pt,
  attach boxed title to top left={xshift=6mm,yshift=-3mm},
  boxed title style={colback=EveFrame,rounded corners},
  coltitle=white,fonttitle=\bfseries\sffamily\small,
  title={\ding{72}\; Everyday picture}}
\newtcolorbox[auto counter]{worked}[1]{enhanced,breakable,colback=WrkTint,colframe=WrkFrame,
  boxrule=0.6pt,arc=4pt,
  attach boxed title to top left={xshift=6mm,yshift=-3mm},
  boxed title style={colback=WrkFrame,rounded corners},
  coltitle=white,fonttitle=\bfseries\sffamily\small,
  title={\ding{46}\; Worked example: #1}}
\newtcolorbox{watchout}{enhanced,breakable,colback=WatTint,colframe=WatFrame,
  boxrule=0.6pt,arc=4pt,
  attach boxed title to top left={xshift=6mm,yshift=-3mm},
  boxed title style={colback=WatFrame,rounded corners},
  coltitle=white,fonttitle=\bfseries\sffamily\small,
  title={\ding{55}\; Watch out}}
\newtcolorbox{keytake}{enhanced,breakable,colback=KeyTint,colframe=KeyFrame,
  boxrule=0.6pt,arc=4pt,
  attach boxed title to top left={xshift=6mm,yshift=-3mm},
  boxed title style={colback=KeyFrame,rounded corners},
  coltitle=white,fonttitle=\bfseries\sffamily\small,
  title={\ding{52}\; Key takeaway}}

% Section style
\titleformat{\section}{\large\bfseries\sffamily\color{SectionNavy}}{}{0em}{}[\titlerule]
\titleformat{\subsection}{\normalsize\bfseries\sffamily\color{SectionNavy}}{}{0em}{}

% Header/footer
\pagestyle{fancy}
\fancyhf{}
\lhead{\small\textsf{ACI Companion}}
\rhead{\small\textsf{Lesson 1 · AI Agents, PEAS \& Environments}}
\renewcommand{\headrulewidth}{0.4pt}
\cfoot{\small\thepage}

\newcommand{\housefig}[2]{\begin{center}\includegraphics[width=\linewidth]{#1}\end{center}}

\begin{document}

% Title banner
\begin{tcolorbox}[enhanced,colback=NavyBanner,colframe=NavyBanner,arc=0pt,
  boxrule=0pt,left=10mm,right=10mm,top=8mm,bottom=8mm]
  {\small\bfseries\sffamily\color{white}\MakeUppercase{ACI · AIMLCZG557}}\par\vspace{4pt}
  {\LARGE\bfseries\sffamily\color{white}AI Agents, PEAS \& Environments}\par\vspace{4pt}
  {\normalsize\color{white!80!NavyBanner}Advanced Computing for AI — Lesson 1}\par
  \textcolor{white!40!NavyBanner}{\rule{\linewidth}{0.4pt}}\par\vspace{2pt}
  {\small\color{white!70!NavyBanner}Exam: 28 Jun 2026 · Closed book · Every slide example worked in full}
\end{tcolorbox}

\vspace{4pt}
{\small\textbf{How to use this companion.} Read the \textcolor{IntFrame}{blue intuition boxes} first — they give the idea in plain words. \textcolor{EveFrame}{Amber everyday boxes} give a real-life analogy. \textcolor{WrkFrame}{Green worked boxes} show a problem solved step by step. \textcolor{WatFrame}{Red watch-out boxes} name the most common mistakes. \textcolor{KeyFrame}{Purple key-takeaway boxes} give the single line to remember.}

\vspace{8pt}

%=============================================================================
\section{1 \quad What is AI?}

Artificial Intelligence is the study of rational agents. A \textbf{rational agent} is anything that perceives its environment and acts to maximise a performance measure. The performance measure is your scoring rule — it says what counts as success.

\begin{intuition}
An agent is not just a robot. It is any decision-maker: a chess program, a spam filter, a thermostat. What makes it \emph{rational} is that it picks the action most likely to score well, given everything it currently knows.
\end{intuition}

\begin{everyday}
Think of a taxi driver. The environment is the city. The sensors are eyes and GPS. The actuators are the steering wheel and pedals. The performance measure is: arrive safely, on time, and cheaply. The driver (agent) picks routes that maximise those goals. An AI taxi agent does the same — but in silicon.
\end{everyday}

The agent--environment loop works like this. At each time step, the agent receives a \textbf{percept} (a reading from its sensors). It processes that percept using its internal mechanism. Then it selects and executes an \textbf{action}. The environment changes. A new percept arrives. The cycle repeats.

\textbf{Where this shows up in ML.} Reinforcement learning is the direct formalisation of this loop. The agent is the policy network; the environment is the simulator; the percept is the observation vector; the action is the output token or motor command; the performance measure is the cumulative reward.

\begin{keytake}
AI = study of rational agents. A rational agent picks the action expected to maximise its performance measure given its percepts so far.
\end{keytake}

%=============================================================================
\section{2 \quad PEAS Specification}

Before building an agent, we describe its task with four components. The acronym is \textbf{PEAS}: \textbf{P}erformance measure, \textbf{E}nvironment, \textbf{A}ctuators, \textbf{S}ensors.

\begin{intuition}
PEAS is a job description for the agent. Performance says what success looks like. Environment says where the agent lives. Actuators say what it can do. Sensors say what it can see.
\end{intuition}

\begin{worked}{Vacuum cleaner PEAS}
\begin{enumerate}[label=\textbf{Step \arabic*.}]
  \item \textbf{Performance measure.} Maximise cleanliness of all squares. Minimise energy used. Minimise time.
  \item \textbf{Environment.} A grid of squares, each either dirty or clean. The agent starts at some square.
  \item \textbf{Actuators.} Move Left, Move Right, Suck (clean the current square).
  \item \textbf{Sensors.} Location sensor (which square am I on?). Dirt sensor (is this square dirty or clean?).
\end{enumerate}
Full PEAS: Performance = clean squares / energy / time; Environment = 2-square grid; Actuators = Left, Right, Suck; Sensors = location, dirt-detector.
\end{worked}

\begin{worked}{Self-driving car PEAS}
\begin{enumerate}[label=\textbf{Step \arabic*.}]
  \item \textbf{Performance.} Safe arrival, minimal travel time, legal driving, passenger comfort.
  \item \textbf{Environment.} Roads, other vehicles, pedestrians, weather, traffic signals.
  \item \textbf{Actuators.} Steering wheel, accelerator, brake, horn, turn signals.
  \item \textbf{Sensors.} Cameras, LIDAR, GPS, speedometer, odometer, microphones.
\end{enumerate}
\end{worked}

\begin{worked}{Sliding-tile adversarial game (exam sample question)}
The exam gives a board with 3 A-tiles, 3 B-tiles, and 2 empty spaces. Players alternate sliding or hopping tiles. The goal is 3 identical tiles adjacent.
\begin{enumerate}[label=\textbf{Step \arabic*.}]
  \item \textbf{Performance.} Win the game (get 3 similar tiles adjacent first).
  \item \textbf{Environment.} The board — a grid with A-tiles, B-tiles, and empty spaces.
  \item \textbf{Actuators.} Slide an adjacent tile into an empty space; hop one tile over another into an empty space.
  \item \textbf{Sensors.} Full view of the current board configuration (fully observable).
\end{enumerate}
\end{worked}

\textbf{Where this shows up in ML.} Every RL problem starts with a PEAS-like specification. In OpenAI Gym, the environment defines the state space (sensors), action space (actuators), and reward function (performance). The ``E'' in PEAS maps directly to the \texttt{env} object.

\begin{keytake}
PEAS = Performance, Environment, Actuators, Sensors. Write all four before designing any agent.
\end{keytake}

%=============================================================================
\section{3 \quad Environment Properties — 6 Pairs}

Every environment is described along six independent dimensions. For exam questions, identify the property and \emph{justify briefly}.

\begin{intuition}
These six pairs are a checklist. You read the problem description, tick each pair, and you have a complete picture of how hard the agent's task will be.
\end{intuition}

\begin{center}
\begin{tabular}{lll}
\toprule
\textbf{Property} & \textbf{One side} & \textbf{Other side} \\
\midrule
Observability & Fully observable & Partially observable \\
Agents & Single agent & Multi-agent (competitive / cooperative) \\
Determinism & Deterministic & Stochastic \\
Episodicity & Episodic & Sequential \\
Dynamics & Static & Dynamic \\
State space & Discrete & Continuous \\
\bottomrule
\end{tabular}
\end{center}

\textbf{Fully observable} means the agent's sensors give it the complete state at every step (chess board). \textbf{Partially observable} means some state is hidden (poker — you can't see opponents' hands).

\textbf{Deterministic} means the next state is completely determined by the current state and the action taken. \textbf{Stochastic} means there is uncertainty (rolling dice, slippery floors).

\textbf{Episodic} means each task is independent — what happened in episode 1 has no effect on episode 2 (sorting mail: each email is separate). \textbf{Sequential} means past actions affect future decisions (chess: every move changes what's possible).

\textbf{Static} means the environment doesn't change while the agent is thinking (crossword puzzle). \textbf{Dynamic} means it can change at any time (driving — other cars keep moving).

\begin{worked}{Classify: online chess against a human}
\begin{enumerate}[label=\textbf{Step \arabic*.}]
  \item Fully observable — both players see the entire board.
  \item Multi-agent (competitive) — two players, each trying to win.
  \item Deterministic — pieces move to exactly the stated square, no dice.
  \item Sequential — every past move restricts future legal moves.
  \item Static — the board does not change while you think (no time pressure changes state).
  \item Discrete — finite number of squares and piece types.
\end{enumerate}
\textbf{Result:} Fully observable, Multi-agent competitive, Deterministic, Sequential, Static, Discrete.
\end{worked}

\begin{watchout}
Do not confuse \emph{static} with \emph{deterministic}. A static environment simply does not change while the agent deliberates. A deterministic environment has no randomness in the outcome of actions. An environment can be deterministic and dynamic at the same time (e.g.\ a fast-moving but physics-exact simulation).

Also: \emph{multi-agent} doesn't always mean \emph{competitive}. Two robots co-operating to move furniture are multi-agent but cooperative.
\end{watchout}

\textbf{Where this shows up in ML.} RL environments in research are often described by exactly these properties. Atari games are fully observable, single-agent, deterministic, episodic, dynamic, discrete. Real-world robotics is partially observable, single-agent, stochastic, sequential, dynamic, continuous — much harder.

\begin{keytake}
Six pairs: (1) observability, (2) number of agents, (3) determinism, (4) episodic vs sequential, (5) static vs dynamic, (6) discrete vs continuous. Memorise all six; the exam asks you to justify each one.
\end{keytake}

%=============================================================================
\section{4 \quad The 5 Agent Types}

Agents are classified by how sophisticated their decision-making is.

\begin{intuition}
Each type adds one more layer of intelligence. Simple reflex just reacts. Model-based remembers the world. Goal-based plans ahead. Utility-based trades off competing goals. Learning agents improve over time.
\end{intuition}

\begin{enumerate}[label=\textbf{\arabic*.}]
  \item \textbf{Simple Reflex Agent.} Condition--action rules only. If dirty, suck. If at edge, turn. No memory of the past. Fails in partially observable environments.
  \item \textbf{Model-Based Reflex Agent.} Keeps an internal model of the world to handle partial observability. The model tracks what it cannot directly see.
  \item \textbf{Goal-Based Agent.} Has an explicit goal. Searches for sequences of actions that reach the goal. More flexible than reflex — can plan.
  \item \textbf{Utility-Based Agent.} Has a utility function mapping states to real numbers. Chooses actions that maximise expected utility. Handles competing goals (fast \emph{and} safe).
  \item \textbf{Learning Agent.} Starts with limited knowledge and improves through experience. Has four components: learning element (updates), performance element (acts), critic (evaluates), problem generator (experiments).
\end{enumerate}

\begin{watchout}
The exam often asks: ``Which agent type is best for environment X?'' A simple reflex agent cannot handle partial observability. A goal-based agent cannot handle competing goals elegantly. Use a utility-based or learning agent for real-world complexity.
\end{watchout}

\textbf{Where this shows up in ML.} Deep RL agents are learning agents: the neural network is the performance element, gradient descent is the learning element, the reward is the critic, and exploration strategies (epsilon-greedy, UCB) are the problem generator.

\begin{keytake}
5 types: Simple reflex → Model-based reflex → Goal-based → Utility-based → Learning agent. Each type is strictly more powerful than the one before it.
\end{keytake}

%=============================================================================
\section{5 \quad Problem-Solving Agents}

A problem-solving agent formulates a search problem and finds a solution. It has four steps: \textbf{Goal formulation} (what do we want?), \textbf{Problem formulation} (what states and actions?), \textbf{Search} (find a sequence of actions), \textbf{Execution} (carry it out).

A formal search problem has five components:
\begin{itemize}
  \item \textbf{State space} — all possible configurations.
  \item \textbf{Initial state} — where we start.
  \item \textbf{Actions} — what we can do.
  \item \textbf{Transition model} — \texttt{Result}$(s, a) \to s'$.
  \item \textbf{Goal test} — is this state the destination?
  \item \textbf{Path cost} — numeric cost of a path (sum of step costs).
\end{itemize}

A \textbf{solution} is a sequence of actions leading from the initial state to a goal state. An \textbf{optimal solution} minimises path cost.

\textbf{Where this shows up in ML.} Hyperparameter search, architecture search (NAS), and planning in model-based RL all reduce to search problems. The state space is the space of hyperparameter combinations or network topologies.

\begin{keytake}
Problem-solving agent = goal formulation + problem formulation + search + execute. A solution is any path to the goal; an optimal solution is the cheapest such path.
\end{keytake}

\end{document}
